\section{Objectives and Contributions}
\label{sec:introduction:objectives}
The central aim of this dissertation is to leverage the expressiveness, computational efficiency, and mathematical rigor of normalizing flows to develop a unified, data-driven framework for probabilistic system identification, state estimation, and optimal control. The primary objectives of this dissertation are as follows:
\begin{itemize}
    \item To investigate the efficacy and limitations of normalizing flows in identifying parameters from system observations.
    \item To adapt and extend normalizing flow architectures for sequential data processing, enabling the estimation of hidden, time-varying states in stochastic dynamical systems.
    \item To integrate and evaluate normalizing flows into control frameworks, exploring their potential to provide uncertainty-aware decision-making within optimal control and reinforcement learning frameworks.
\end{itemize}

Building upon these objectives, the key contributions of this dissertation are threefold:
\begin{itemize}
    \item Parameter Identification
    \begin{itemize}
        \item Unification of a Deep-Generative-Model-Based Bayesian Inference Framework: A comprehensive framework that leverages deep generative models for parameter identification, comparing their performance and characteristics across different scenarios.
        \item Investigating the Relationship Between Sensitivity and Identifiability: Integration of sensitivity analysis with normalizing flows to enable reliable parameter identification in systems where parameters exhibit varying levels of identifiability.
        \item Combination of MCMC with Normalizing Flows: The integration of Markov Chain Monte Carlo methods with normalizing flows to achieve robust parameter identification under observational and system uncertainty.
        \item Real-World Application of Normalizing Flows for Parameter Identification: The application of normalizing flows to real-world systems, demonstrating their practical utility and effectiveness in parameter identification tasks.
    \end{itemize} 
    \item State Estimation: Development of novel state estimation methods, including amortized sequence flow, latent state-space model flow, and conditional stochastic interpolant, to capture complex, non-Gaussian uncertain state trajectories from observations in stochastic dynamical systems.
    \item Optimal Control: Integration of flow-based approaches into stochastic control frameworks to enable uncertainty-aware decision-making in complex, unpredictable environments.
\end{itemize}